\section{Appendix}
\label{app:top}

\subsection{Models and datasets}
\label{app:models}
Table~\ref{tab:models} lists the explained models. Training uses a stratified 70/30 split with a fixed seed; the models are not tuned. XGBoost: depth 4, 200 trees, learning rate 0.1, subsample 0.8; random forest: 100 trees, depth at most 8, minimum leaf size 3; MLP: two hidden layers of 64 units, $L_2$ penalty $10^{-2}$; LR: $C{=}1$. The 50-row background is drawn from the training rows with a fixed seed. The 50 instances per dataset are 25 random test rows and the 25 of a random pool of 300 test rows nearest to the decision boundary.

\begin{table}[h]
\centering
\caption{Explained models, the value function explained, and held-out performance (one fit per model, seed 0).}
\label{tab:models}
\input{../tables/tab5_models.tex}
\end{table}

\subsection{Controls, per-instance distributions}
\label{app:controls}
Figure~\ref{fig:instances} shows, for each explained instance on the stochastic datasets, the mean top-3 Jaccard between two runs at $\nicefrac14\times$ and $1\times$. Figure~\ref{fig:controls} repeats the budget curves on the two deterministic-regime controls, where the library default enumerates all coalitions and the budget curves are flat by construction.

\begin{figure}[h]
\centering
\includegraphics[width=\linewidth]{fig2_instance_distribution.pdf}
\caption{Per-instance mean pairwise top-3 Jaccard between runs (190 run pairs per instance) at $\nicefrac14\times$ and $1\times$ budget, for every explained instance, coloured by stratum (random or near boundary); black bar: mean. A minority of instances lie well below the mean in both strata.}
\label{fig:instances}
\end{figure}

\begin{figure}[h]
\centering
\includegraphics[width=\linewidth]{fig1b_budget_curves_controls.pdf}
\caption{As Figure~\ref{fig:budget} for the controls ACS Income ($d{=}10$) and Pima diabetes ($d{=}8$), where coalition enumeration is used once the budget allows it.}
\label{fig:controls}
\end{figure}

\subsection{Per-model verdicts and cells not run}
\label{app:cells}
Table~\ref{tab:claims} gives, for every cell, the metrics entering the decision rule; Table~\ref{tab:notrun} lists the cells not run and why. Not run at all: the background-size axis, the Give Me Some Credit and Arrhythmia datasets, and comparisons with other explainers; they were cut for compute budget before the runs started and nothing was run and discarded, except that partially finished $16\times$ ACS Employment cell files were deleted.

{\footnotesize
\input{../tables/tab6_claims.tex}
\addtocounter{table}{-1}% longtable has advanced the counter
\captionof{table}{Per cell: budget $B$ (model evaluations), instances/runs, top-3 hit rate against the exact target with its 95\% interval, median relative $L_2$ error, lower bound of top-3 sign agreement, top-1 Jaccard between runs, and the verdict.}
\label{tab:claims}
}

\input{../tables/tab6b_cells_not_run.tex}
\addtocounter{table}{-1}% longtable has advanced the counter
\captionof{table}{Cells not run, with reasons.}
\label{tab:notrun}

\subsection{Sensitivity of the verdicts to the thresholds}
\label{app:sens}
Table~\ref{tab:verdictsens} recounts the cells (reliable / conditional / unreliable) when the thresholds of the rule are varied one at a time and together, on German Credit and ACS Employment, over all arms including the $\ell_1$-default arm. The counts are computed from the bounds already in \texttt{verdicts.csv}; the pre-stated row reproduces the stored verdicts, and the cells are not independent (the same instances serve every budget and arm). The additive LR is exact in every cell and the tree models are explained with an exact target, so only the MLP column is evidence about a setting where sampling is the practical option. With $\ell_1$ off, the number of reliable MLP cells falls from $\nSensMlpBase$ of $\nSensMlpN$ under the pre-stated rule to $\nSensMlpStrictA$ and $\nSensMlpStrictB$ under the two strict sets; on German Credit all six MLP cells fail the strictest set (lowest hit-rate lower bound $\nSensGerMlpHitLoMin$, largest median relative error $\nSensGerMlpRelErrMax$). The verdict ``reliable'' is therefore threshold-sensitive for the one informative model.

\begin{table}[h]
\centering
\caption{Cells classed reliable / conditional / unreliable under alternative threshold sets $(h,e,s)$ (top-3 hit lower bound, median relative $L_2$, sign-agreement lower bound; the criteria for ``unreliable'' are unchanged), for all models, without LR, and for the MLP alone. The number of cells is in the header.}
\label{tab:verdictsens}
\input{../tables/tab11_verdict_sensitivity.tex}
\end{table}

\subsection{Stricter notions, strata and tails}
\label{app:strata}
Table~\ref{tab:stricter} reports, at $1\times$ with $\ell_1$ off, stricter agreement notions than the top-3 set: top-1, top-3 and top-5 hit, the top-3 set in the target's order, and the sign of all features. Ordering is the weakest notion: the ordered top-3 hit is $0.927$ for the German MLP (20 instances) against $0.977$ for the unordered set. Tables~\ref{tab:strata} and~\ref{tab:tail} split the 50 instances into the two sampling strata (25 random, 25 near the boundary) and give the share of instances whose per-instance mean top-3 Jaccard falls below $0.99$ to $0.80$. The explained set is a balanced mix by design, so neither pooled value is a population fraction, and with 25 instances per stratum the intervals are wide; we draw no conclusion about which stratum is harder. Not computed: the per-(instance $\times$ run) top-3 hit fraction and the top-10 hit.

\begin{table}[h]
\centering
\caption{Stricter notions at $1\times$ (Kernel SHAP, $\ell_1$ off): mean with nested-bootstrap 95\% interval; $n$ instances (German MLP: 20).}
\label{tab:stricter}
\input{../tables/tab12_stricter.tex}
\end{table}

{\footnotesize
\input{../tables/tab13_strata.tex}
\addtocounter{table}{-1}% longtable has advanced the counter
\captionof{table}{Main $1\times$ metrics per instance stratum (random versus near the decision boundary), instance-bootstrap 95\% intervals within each stratum (1000 draws); German MLP error rows rest on fewer instances.}
\label{tab:strata}
}

{\footnotesize
\input{../tables/tab13b_tail.tex}
\addtocounter{table}{-1}% longtable has advanced the counter
\captionof{table}{Percentage of instances whose per-instance mean top-3 Jaccard between runs is below $0.99$, $0.95$, $0.90$ and $0.80$, per stratum, with instance-bootstrap intervals.}
\label{tab:tail}
}

\subsection{Budget contrasts, equal-budget comparison and N4 ratios}
\label{app:contrasts}
Tables~\ref{tab:eqbud}--\ref{tab:n4ratio} give, for every cell, the paired contrasts between budgets, the equal-budget averaging comparison and the N4 variance ratios with their intervals. Intervals are nested-bootstrap percentile intervals (instances, then runs); Holm $p$-values correct over the contrasts of one family. Rows of exact or additive models (LR, and tree models on the exact enumeration cells) are included for completeness; their contrasts are at floating-point level and carry no evidence about sampling noise.

\input{../tables/tab8_equal_budget.tex}
\addtocounter{table}{-1}% longtable has advanced the counter
\captionof{table}{Equal-budget averaging: difference (averaged $-$ single run) in top-3 hit rate and relative $L_2$ error at the same total budget. ``$4{\times}\tfrac14\times$'' is the mean of four runs at $\tfrac14\times$ compared with one run at $1\times$; ``$4{\times}1\times$'' is four runs at $1\times$ against one at $4\times$. A positive $\Delta$ rel.\ $L_2$ means the average of runs is worse.}
\label{tab:eqbud}

\input{../tables/tab9_contrasts.tex}
\addtocounter{table}{-1}% longtable has advanced the counter
\captionof{table}{Planned paired contrasts per cell (A vs B; $\Delta$ is A minus B): budget steps, permutation versus Kernel SHAP at the same budget, and $\ell_1$ off versus the library default at $1\times$. Holm-adjusted $p$ for the relative-$L_2$ contrast.}
\label{tab:contrasts}

\input{../tables/tab10_n4_ratios.tex}
\addtocounter{table}{-1}% longtable has advanced the counter
\captionof{table}{N4 probe: explainer-run SD, SD induced by redrawing the background ($L$ levels $\times$ $S$ seeds) and the ratio of the two variances, with bootstrap intervals. The ratio is descriptive for $m=50$ at the $1\times$ budget and is undefined where the explainer is exact.}
\label{tab:n4ratio}

\subsection{Value function}
\label{app:valfn}
Table~\ref{tab:valfn} compares the explained scalar: log-odds against predicted probability, same models, instances, background and seeds (ACS Employment, $1\times$). Reliability stays high with the probability (top-3 hit $\nVfProbTopHitMin$--$\nVfProbTopHitMax$ against $\nVfLogTopHitMin$--$\nVfLogTopHitMax$; Spearman $\nVfProbSpearMin$--$\nVfProbSpearMax$ against $\nVfLogSpearMin$--$\nVfLogSpearMax$), with a somewhat larger relative error ($\nVfProbRelErrMin$--$\nVfProbRelErrMax$ against about $\nVfLogRelErrMin$), as expected, since the probability is a nonlinear transform of the log-odds, and the additive LR is no longer exactly additive in this scale.

\begin{table}[h]
\centering
\caption{Value-function sensitivity on ACS Employment at $1\times$: log-odds versus probability, mean over 50 instances with 95\% bootstrap intervals.}
\label{tab:valfn}
\input{../tables/tab7_valfn.tex}
\end{table}

\subsection{Use of LLMs and compute}
\label{app:llm}
The study (design, implementation, experiments, analysis and writing of this report) was carried out by LLM-based agents (Claude models) inside an automated research environment, working from a written task specification, with separate agents responsible for planning, experiments, figures, venue formatting and review. All numbers in the text come from result files produced by the released code; they enter the text through a generated macro file. Total compute was about 26 CPU-hours on a shared, oversubscribed CPU host (no GPU).
